%% ================================================================================================
%%  9. Conclusion
%% ================================================================================================
\section{Conclusion}\label{sec:conclusion}

The research consensus on veto rights is negative about unanimity and positive about vetoes that are
separate from approval and limited to blocking what is unacceptable. Social choice adds that no procedure
escapes manipulation, that the majority principle and a minority veto must be traded against each other,
and that an open agenda requires independence of clones. These findings point to one design: a single
ballot on which every voter ranks every option and marks the unacceptable ones; the exclusion of the
options rejected by a blocking minority; a guaranteed decision when every option is rejected; and the
Schulze winner among the options left. With a blocking third, this procedure, \QVC(1/3), never deadlocks,
elects an option acceptable to more than two thirds of the voters whenever there is one, respects the
majority among such options, ignores flooding, and in a polarised electorate halves the share of
contested outcomes compared with the Schulze method, at no cost in total welfare. Its weakness is the
weakness of any veto: a blocking minority can use its power strategically. That risk is contained by
procedure, with named and reasoned vetoes, a quota no lower than a quarter and at most one revision
round, and it is the price of the protection the veto gives.
